\paragraph{Exact local image at the coefficient primes.}
Let $p\mid n$, put $s=v_p(n)$, and assume $s\ge3$ when
$p=2$.  Set
\[
 z_0=1+p^s,\qquad
 \mathcal B=z_0+p^{s+1}\Z_p,\qquad
 F_n(z)=\exp\bigl(24n(z^2-1)\log z\bigr)
                   \quad(z\in1+p^s\Z_p).
\tag{A.59.C1}
\]
This definition fixes the Teichmuller component to one.
For positive rational integers $z$ in this ball it equals
the ordinary integer power $z^{24n(z^2-1)}$.
Put $h=3s+v_p(48)$ and
\[
 \epsilon_p=\begin{cases}1,&p\ne2,\\3,&p=2.\end{cases}
\]
Then the exact image is
\[
                 F_n(\mathcal B)
                   =F_n(z_0)+p^{h+\epsilon_p}\Z_p.
\tag{A.59.C2}
\]
For any $a,b,c\in\Z_p$ of which at least one is a unit,
the image of the three-variable map on $\mathcal B^3$ is
\[
 \{aF_n(q)+bF_n(r)-cF_n(t):q,r,t\in\mathcal B\}
       =(a+b-c)F_n(z_0)+p^{h+\epsilon_p}\Z_p.
\tag{A.59.C3}
\]
Consequently this expression can be made zero, with all
three variables retaining exact valuation
$v_p(q-1)=v_p(r-1)=v_p(t-1)=s$, exactly when
\[
                         v_p(a+b-c)\ge h+\epsilon_p.
\tag{A.59.C4}
\]
One may leave two variables equal to $z_0$ and vary only
a variable with unit coefficient.

\emph{Proof of the local image formula.}
We work in the ring $R=\Z_p\langle w\rangle$ of restricted
power series, whose coefficients tend to zero $p$-adically.
Write $z=1+p^s w$ and set
\[
 A=48np^{2s},\qquad v_p(A)=h.
\]
Expansion of the logarithm gives the exact identity
\[
 \begin{split}
 24n\bigl((1+p^s w)^2-1\bigr)\log(1+p^s w)
      &=A w^2V(w),\\
 V(w)&=1+\sum_{k\ge2}
       \frac{(-1)^k(k-1)p^{sk}}{2k(k+1)}w^k.
 \end{split}
\tag{A.59.C5}
\]
There is no linear term in $V$.  To check the displayed
coefficient, multiply
$\log(1+x)/x=\sum_{k\ge0}(-1)^kx^k/(k+1)$ by
$1+x/2$.  The coefficient of $x^k$ for $k\ge1$ is
$(-1)^k(k-1)/(2k(k+1))$.

For odd $p$, the coefficients of $V-1$ belong to $pR$.
Indeed, for $k\ge2$,
$v_p(k(k+1))\le k-1$, so their valuations are at least
$sk-(k-1)\ge1$.  They tend to infinity since the
denominator valuations grow only logarithmically.
At two the same estimate gives
\[
      v_2\left(\frac{(k-1)2^{sk}}{2k(k+1)}\right)
          \ge k(s-1)\ge2s-2\ge4.
\]
Thus
\[
 V\equiv1\pmod{pR}\quad(p\ne2),\qquad
 V\equiv1\pmod{16R}\quad(p=2).
\tag{A.59.C6}
\]
The assumption $s\ge3$ at two also guarantees that all
logarithms and exponentials in (A.59.C5) converge on this
entire ball.

The exponential series in the complete ring $R$ now
has the form
\[
 F_n(1+p^s w)=1+A f(w),\qquad
 f(w)=w^2V(w)+
       \sum_{j\ge2}\frac{A^{j-1}}{j!}(w^2V(w))^j.
\tag{A.59.C7}
\]
This is a restricted integral series.  For example
$v_p(j!)\le j-1$ shows both integrality and convergence
of the second sum, since $h\ge3$ for odd $p$ and
$h\ge13$ at two.  That sum is in $pR$ for odd $p$
and in $16R$ at two.  Hence
\[
 f(w)=w^2+pR_1(w)\quad(p\ne2),\qquad
 f(w)=w^2+16R_1(w)\quad(p=2)
\tag{A.59.C8}
\]
for an integral restricted series $R_1$.

An integral restricted series is one-Lipschitz on
$\Z_p$, and its difference under $w\mapsto1+p^r t$
is divisible by $p^r$ in $\Z_p\langle t\rangle$.
For odd $p$, (A.59.C8) therefore gives
\[
       \frac{f(1+pt)-f(1)}p=2t+pR_2(t),
                   \qquad R_2\in\Z_p\langle t\rangle.
\tag{A.59.C9}
\]
The right side is a bijection of $\Z_p$: for a specified
value $y$, the equation
$t=(y-pR_2(t))/2$ has a unique solution by contraction,
because two is a unit and $R_2$ is one-Lipschitz.
It follows that
$f(1+p\Z_p)=f(1)+p\Z_p$.

At two, every odd square is one modulo eight, so
(A.59.C8) first gives
\[
                 f(1+2\Z_2)\subset f(1)+8\Z_2.
\]
On the smaller ball $w=1+4t$, the same equation gives
\[
        \frac{f(1+4t)-f(1)}8=t+2R_3(t),
                    \qquad R_3\in\Z_2\langle t\rangle.
\tag{A.59.C10}
\]
To see the divisibility of the error directly, the square
term contributes $t+2t^2$, while
$16(R_1(1+4t)-R_1(1))/8$ is divisible by eight.
The right side of (A.59.C10) is again a bijection of
$\Z_2$, by the contraction equation
$t=y-2R_3(t)$.  Thus the smaller ball already maps onto
$f(1)+8\Z_2$, proving equality for $1+2\Z_2$ as well.

Since $z\in\mathcal B$ corresponds exactly to
$w\in1+p\Z_p$, multiplying these exact images by $A$
in (A.59.C7) proves (A.59.C2).  All changes in the
three-variable expression consequently lie in the
ideal on the right of (A.59.C3).  A variable with unit
coefficient fills that ideal by itself while the
other two variables stay at $z_0$.
This proves (A.59.C3), and (A.59.C4) follows because
$F_n(z_0)$ is a unit.  Membership in $\mathcal B$
retains the required exact valuation $s$ for each
variable.

\paragraph{Precision and the derivative check.}
The derivative with respect to $z$ at $z_0$ is
\[
 \frac{F_n'(z_0)}{F_n(z_0)}
   =24n\left(2z_0\log z_0+\frac{z_0^2-1}{z_0}\right),
 \qquad v_p(F_n'(z_0))=2s+v_p(96).
\tag{A.59.C11}
\]
Indeed the bracket is
$4p^s+(2/3)p^{3s}-(5/6)p^{4s}+\cdots$;
the remainder has strictly larger valuation than
the leading term under the assumptions above.
Equivalently (A.59.C7)--(A.59.C8) give the same valuation
after multiplying the derivative in $w$ by $p^{-s}$.
For the invariant proposed in the recursion,
$H_j(p)=2(2s+v_p(96))+j+3$,
one has $H_j(p)\ge h+\epsilon_p$ at every coefficient
prime.  Thus that invariant suffices for (A.59.C4);
the stronger bound using the derivative in $w$
also suffices.

Finally, approximate any chosen local solution by rational
integers modulo $p^L$, with $L\ge s+1$.  The variables
remain in $\mathcal B$.  Their $w$-coordinates change by
elements of $p^{L-s}\Z_p$; the one-Lipschitz property
of $f$ in (A.59.C7) makes each $F_n$ value change by an
element of $p^{h+L-s}\Z_p$.
Taking $L$ arbitrarily large therefore imposes any
specified next-stage error precision.  At the integer
prime variables supplied by the later CRT and
Chebotarev step, (A.59.C1) is the original integer power,
so this precision statement applies to the actual
prime recursion.
